\documentclass[10pt,a4paper]{amsart}
\usepackage[margin=19mm]{geometry}
\usepackage{amsmath,amssymb,mathtools}
\usepackage[scr=rsfs,cal=euler]{mathalfa}
\usepackage{xfrac}
\usepackage[unicode,hidelinks,bookmarks=false]{hyperref}
\newcommand{\ie}{i.e.\ }
\newcommand{\cf}{cf.\ }
\newcommand{\resp}{respectively\ }
\newcommand{\Subsection}{\subsection}
\providecommand{\nobreakdash}{\nobreak\hbox{-}}
\setlength{\emergencystretch}{2em}
\multlinegap0pt
\usepackage{amssymb}
\usepackage{float}
\usepackage{algorithm}\usepackage{algpseudocode}
\newtheorem{lemma}{Lemma}
\newtheorem{definition}{Definition}
\newtheorem{theorem}{Theorem}
	\DeclareMathOperator{\Div}{div}
\def\Om{\Omega}
\def\crr{\cr\noalign{\vskip-0.0mm}}
\def\pt{\partial}
\def\vp{\varphi}
\def\wh{\widehat}
\title{\bf {\bf Weak Unique Continuation on Annular Domains for Backward Degenerate Parabolic Equations with Degenerate Interior Points }}
\author{Dong-Hui Yang, Bao-Zhu Guo, Guojie Zheng, Jie Zhong}
\date{Source: arXiv:2605.02797v1 (CC BY 4.0)}
\begin{document}
\maketitle

\noindent\emph{Source and licence.} \bf {\bf Weak Unique Continuation on Annular Domains for Backward Degenerate Parabolic Equations with Degenerate Interior Points }. Version arXiv:2605.02797v1 (CC BY 4.0), distributed by arXiv under the Creative Commons Attribution 4.0 International licence, http://creativecommons.org/licenses/by/4.0/, as recorded in the arXiv licence metadata for this exact version and shown on the abstract page https://arxiv.org/abs/2605.02797v1. Full licence notice: \texttt{licenses/arxiv-2605.02797-CC-BY-4.0.txt}.\par\smallskip

\noindent\textit{Editorial scope.} Counted unit: the article's theorem 08.16.T1 (source0.tex lines 867--879). The article's complete original proof of it is delivered with it (source0.tex lines 880--988, one \texttt{proof} environment of 109 lines). No mathematics is written, repaired or re-derived: the statement and the proof are the article's own lines, in the article's own notation, and no retained line is substituted or re-flowed.  Retained uncounted as the source-local context the proof needs: lines 490--496; lines 580--586; lines 660--674; lines 818--824; lines 841--851.  Nothing was omitted from any retained span, so the package carries no omission record.  No retained line contains a citation, so the wrapper carries no bibliography and no bibliography entry.  No retained line contains a figure environment, an image-inclusion command or drawing code, so no image is delivered and the package declares no asset dependency.  Editorial formatting only, in addition: an AMS \texttt{amsart} preamble that restates the article's own declarations and macros used by the retained lines, namely the environments \texttt{lemma}, \texttt{definition}, \texttt{theorem} on separate counters, together with the standard packages those lines need.  The retained environments print in this document as Lemma 1, Definition 1, Theorem 1 in order of appearance, whereas the article prints the same items with the numbering of its own full document.  The complete native source of the article as distributed in the arXiv e-print for this exact version is bundled unchanged as \texttt{sources/199749/source0.tex}.
\begin{equation}\label{12.14.1}
\begin{cases}
\partial_t \varphi + \Div(|x|^\alpha \nabla \varphi) = 0, & \text{ in } Q, \\
\varphi = 0, & \text{ on } \partial Q, \\
\varphi(T) = \varphi_T, & \text{ in } \Omega,
\end{cases}
\end{equation}

\begin{lemma}\label{08.15.L1}
For every $N \geq 2$ and $\alpha \in (0, 2)$, it holds that
\begin{equation*}
(N - 2 + \alpha)\left\||x|^{\frac{\alpha}{2} - 1}u\right\|_{L^2(\Omega)} \leq 2\|\nabla u\|_{L^2(\Omega; w)}
\end{equation*}
for all $u \in H_0^1(\Omega; w)$. Furthermore, if $u \in H_0^1(\Omega; w)$, then $u \in L^2(\Omega)$.
\end{lemma}

\begin{definition}\label{08.16.D1}
Let $f \in L^2(Q; w)$ and $\varphi_0 \in L^2(\Omega)$. We call $\varphi \in W$ a \emph{weak solution} of
\begin{equation}\label{03.01.1}
\begin{cases}
\partial_t \varphi - \mathrm{Div}(|x|^\alpha \nabla \varphi) = f, &\text{ in } Q, \\
\varphi = 0, &\text{ on } \partial Q, \\
\varphi(0) = \varphi_0, &\text{ in } \Omega,
\end{cases}
\end{equation}
if
\begin{equation*}
-\iint_Q \varphi \partial_t \psi \, dx \, dt + \iint_Q (\nabla \varphi \cdot \nabla \psi)w \, dx \, dt = \iint_Q f \psi \, dx \, dt + \int_\Omega \varphi_0(x)\psi(x, 0) \, dx
\end{equation*}
for any $\psi \in W$ with $\psi(T) = 0$.
\end{definition}

\begin{equation}\label{08.16.1}
\begin{cases}
\pt_t\wh \vp_k - \Div(w_k\nabla\wh \vp_k) = f_k, &\text{ in } Q, \\
\wh \vp_k = 0, &\text{ on } \pt Q, \\
\wh \vp_k(0) = \vp_k, & \text{ in } \Om,
\end{cases}
\end{equation}

  \begin{lemma}\label{08.22.L1}
Let $\varphi_k \in W_k$ ($k \in \mathbb{N}$) denote the solution of \eqref{08.16.1}, where $\varphi_0 \in L^2(\Omega)$ and $f_k w_k^{-1} \in L^2(\Omega; w_k)$. Then, the following estimate holds:
\begin{equation*}
\max_{t \in [0,T]} \|\varphi_k(t)\|_{L^2(\Omega)} + \|\varphi_k\|_{L^2(0,T; H_0^1(\Omega; w_k))} \leq C \left( \|\varphi_0\|_{L^2(\Omega)} + \|f_k w_k^{-1}\|_{L^2(\Omega; w_k)} \right),
\end{equation*}
where the positive constant $C$ depends solely on $\alpha$, $N$, and $\Omega$.
Moreover, if $f_k = 0$, then for all $0 \leq t \leq \tau \leq T$, the following inequality is satisfied:
\begin{equation*}
\|\varphi_k(\tau)\|_{L^2(\Omega)} \leq \|\varphi_k(t)\|_{L^2(\Omega)}.
\end{equation*}
\end{lemma}

 \begin{theorem}\label{08.16.T1}
Let $\widehat{\varphi}_0 \in W$ denote the solution of equation \eqref{12.14.1} with initial data $\varphi_0 \in L^2(\Omega)$ and $fw^{-1} \in L^2(\Omega;w)$. For $k \in \mathbb{N}$, let $\widehat{\varphi}_k$ be the solution of equation \eqref{08.16.1} with initial data $\varphi_k = \varphi_0$ and source term $f_k = f$. Then, the following weak convergence results hold:
\begin{equation}\label{08.22.1}
\begin{split}
\widehat{\varphi}_k &\rightharpoonup \widehat{\varphi}_0   \text{  weakly in } L^2(0,T; H_0^1(\Omega;w)), \\
\widehat{\varphi}_k &\rightharpoonup \widehat{\varphi}_0   \text{  weakly in } L^2(Q).
\end{split}
\end{equation}
Furthermore, we have:
\begin{equation}\label{12.09.5}
\widehat{\varphi}_k(T) \rightharpoonup \widehat{\varphi}_0(T) \text{  weakly in } L^2(\Omega).
\end{equation}
\end{theorem}

\begin{proof}  The proof will be divided into several steps.

{\it Step 1.}  Observe that
\begin{equation}\label{12.11.1}
\begin{split}
\iint_Q (fw_k^{-1})^2w_k\,\mathrm{d} x\,\mathrm{d} t &= \iint_{Q}(fw^{-1})^2(ww_k^{-1})w\,\mathrm{d} x\,\mathrm{d} t \\
&\leq \iint_Q (fw^{-1})^2w\,\mathrm{d} x\,\mathrm{d} t,
\end{split}
\end{equation}
and
\begin{equation*}
\iint_{\Omega\times (0,t)}|\nabla\widehat{\varphi}_k|^2w\,\mathrm{d} x\,\mathrm{d} t \leq \iint_{\Omega\times (0,t)}|\nabla \widehat{\varphi}_k|^2w_k\,\mathrm{d} x\,\mathrm{d} t,
\end{equation*}
since $w\leq w_k$ in $\Omega$. By Lemma \ref{08.22.L1}, we have
\begin{equation*}
\begin{split}
\max_{t\in [0,T]}\|\widehat{\varphi}_k(t)\|_{L^2(\Omega)}+\|\widehat{\varphi}_k\|_{L^2(0,T; H_0^1(\Omega;w))} \leq C\left(\|\varphi_0\|_{L^2(\Omega)}+\|fw^{-1}\|_{L^2(\Omega; w)}\right),
\end{split}
\end{equation*}
where the constant $C>0$ depends only on $\alpha$, $N$, and $\Omega$. Consequently, there exists a subsequence of $\{\widehat{\varphi}k\}{k\in\mathbb{N}}$, still denoted by itself, and $\widetilde{\varphi}_0\in L^2(0,T; H_0^1(\Omega;w))$, such that
\begin{equation}\label{08.16.4}
\begin{split}
\widehat{\varphi}_k
&\rightharpoonup \widetilde{\varphi}_0   \text{  weakly in } L^2(0,T; H_0^1(\Omega;w)),
\end{split}
\end{equation}
and by Lemma \ref{08.15.L1}, we obtain
\begin{equation}\label{08.23.1}
\widehat{\varphi}_k
\rightharpoonup \widetilde{\varphi}_0   \text{  weakly in } L^2(Q).
\end{equation}
Furthermore, we have
\begin{equation*}
\widehat{\varphi}_k
\rightharpoonup \widetilde{\varphi}_0   \text{  weakly in } L^2(Q; w).
\end{equation*}

 {\it Step 2}. Next, we aim to prove that $\widetilde{\varphi}_0$ constitutes a weak solution of \eqref{12.14.1}.
Assume $\psi\in W$ with $\psi(T)=0$. By leveraging a density argument, we can assume $\psi\in C^\infty(\overline{Q})$ and that $\psi(\cdot, t)$ is compactly supported in $\Omega$ for any $t$. Utilizing the definitions of weak solution (refer to Definitions \ref{08.16.D1} and \ref{08.16.D1}) and equations \eqref{08.16.4} and \eqref{08.23.1}, it suffices to demonstrate
\begin{equation}\label{08.16.5}
\iint_Q(\nabla\widehat{\varphi}_k\cdot \nabla\psi)w_k\mathrm{d} x\mathrm{d} t\to \iint_Q(\nabla\widetilde{\varphi}_0\cdot\nabla\psi)w\mathrm{d} x\mathrm{d} t   \text{  as } k\to \infty.
\end{equation}
Firstly, for any $\gamma>0$, by \eqref{08.16.4}, there exists $k_0\in\mathbb{N}$ such that for all $k\geq k_0$, we obtain
\begin{equation}\label{08.16.6}
\left|\iint_Q(\nabla\widehat{\varphi}_k\cdot\nabla\psi)w\mathrm{d} x\mathrm{d} t-\iint_Q(\nabla\widetilde{\varphi}_0\cdot\nabla\psi)w\mathrm{d} x\mathrm{d} t\right|<\frac{1}{2}\gamma.
\end{equation}
Secondly, observe that $w=w_k$ on $\Omega\setminus B_{\frac{1}{k}}$, implying
\begin{equation}\label{08.16.7}
\iint_{(\Omega\setminus B_{\frac{1}{k}})\times (0,T)}(\nabla\widehat{\varphi}_k\cdot\nabla\psi)w_k\mathrm{d} x\mathrm{d} t=\iint_{(\Omega\setminus B_{\frac{1}{k}})\times (0,T)}(\nabla \widehat{\varphi}_k\cdot\nabla\psi)w\mathrm{d} x\mathrm{d} t.
\end{equation}
 Thirdly, by utilizing the same reasoning as employed in \eqref{12.11.1} from Step 1, we obtain
\begin{eqnarray}\label{08.16.9}
&&\left|\iint_{B_{\frac{1}{k}} \times (0,T)} (\nabla\widehat{\varphi}_k \cdot \nabla\psi) w_k \, dx \, dt\right| \crr
&&\leq \left(\iint_{B_{\frac{1}{k}} \times (0,T)} |\nabla\widehat{\varphi}_k|^2 w_k \, dx \, dt\right)^{\frac{1}{2}} \left(\iint_{B_{\frac{1}{k}} \times (0,T)} |\nabla\psi|^2 w_k \, dx \, dt\right)^{\frac{1}{2}} \crr
&&\leq C\sqrt{T} \left(\sup_{Q} |\nabla\psi|\right) \left(\|\varphi_0\|_{L^2(\Omega)} + \|f w^{-1}\|_{L^2(Q; w)}\right) \left(w_k(B_{\frac{1}{k}})\right)^{\frac{1}{2}} \crr
&&\leq C_\psi \sqrt{T} \left(\|\varphi_0\|_{L^2(\Omega)} + \|f w^{-1}\|_{L^2(\Omega; w)}\right) \frac{1}{k^{\frac{N+\alpha}{2}}}
\end{eqnarray}
given that
\begin{equation*}
w_k(B_{\frac{1}{k}}) = \int_{B_{\frac{1}{k}}} w_k \, dx = \int_{B_{\frac{1}{k}}} \frac{1}{k^\alpha} \, dx = C \frac{1}{k^{N+\alpha}},
\end{equation*}
where the constant $C_\psi$ depends exclusively on $\alpha, N, \Omega,$ and $\psi$.
Fourthly,
\begin{equation}\label{08.16.10}
\begin{split}
&\left|\iint_{B_{\frac{1}{k}} \times (0,T)} (\nabla\widehat{\varphi}_k \cdot \nabla\psi) w \, dx \, dt\right| \\
&\leq \left(\iint_{B_{\frac{1}{k}} \times (0,T)} |\nabla\widehat{\varphi}_k|^2 w \, dx \, dt\right)^{\frac{1}{2}} \left(\iint_{B_{\frac{1}{k}} \times (0,T)} |\nabla\psi|^2 w \, dx \, dt\right)^{\frac{1}{2}} \\
&\leq C\sqrt{T} \left(\sup_{Q} |\nabla\psi|\right) \left(\|\varphi_0\|_{L^2(\Omega)} + \|f\|_{L^2(Q; w)}\right) \left(w(B_{\frac{1}{k}})\right)^{\frac{1}{2}} \\
&\leq C_\psi \sqrt{T} \left(\|\varphi_0\|_{L^2(\Omega)} + \|f\|_{L^2(Q; w)}\right) \frac{1}{k^{\frac{N+\alpha}{2}}}
\end{split}
\end{equation}
in accordance with
\begin{equation*}
w(B_{\frac{1}{k}}) = \int_{B_{\frac{1}{k}}} w \, dx = \int_{B_{\frac{1}{k}}} |x|^\alpha \, dx = C \frac{1}{k^{N+\alpha}},
\end{equation*}
where the constant $C_\psi$ is contingent only on $\alpha, N, \Omega,$ and $\psi$.
Ultimately, since $w = w_k$ on $\Omega \setminus B_{\frac{1}{k}}$, by \eqref{08.16.6}, \eqref{08.16.7}, \eqref{08.16.9}, and \eqref{08.16.10}, we derive
\begin{eqnarray*}
&&\left|\iint_Q (\nabla\widehat{\varphi}_k \cdot \nabla\psi) w_k \, dx \, dt - \iint_Q (\nabla\widetilde{\varphi}_0 \cdot \nabla\psi) w \, dx \, dt\right| \\
&&\leq \left|\iint_Q (\nabla\widehat{\varphi}_k \cdot \nabla\psi) w_k \, dx \, dt - \iint_Q (\nabla\widehat{\varphi}_k \cdot \nabla\psi) w \, dx \, dt\right| \\
&&  + \left|\iint_Q (\nabla\widehat{\varphi}_k \cdot \nabla\psi) w \, dx \, dt - \iint_Q (\nabla\widetilde{\varphi}_0 \cdot \nabla\psi) w \, dx \, dt\right| \\
&&\leq \left|\iint_{(\Omega \setminus B_{\frac{1}{k}}) \times (0,T)} (\nabla\widehat{\varphi}_k \cdot \nabla\psi) w_k \, dx \, dt - \iint_{(\Omega \setminus B_{\frac{1}{k}}) \times (0,T)} (\nabla\widehat{\varphi}_k \cdot \nabla\psi) w \, dx \, dt\right| \\
&&  + \left|\iint_{B_{\frac{1}{k}} \times (0,T)} (\nabla\widehat{\varphi}_k \cdot \nabla\psi) w_k \, dx \, dt\right| + \left|\iint_{B_{\frac{1}{k}} \times (0,T)} (\nabla\widehat{\varphi}_k \cdot \nabla\psi) w \, dx \, dt\right| + \frac{1}{2}\gamma \\
&&\leq C_\psi \sqrt{T} \left(\|\varphi_0\|_{L^2(\Omega)} + \|f\|_{L^2(Q; w)}\right) \frac{1}{k^{\frac{N+\alpha}{2}}} + \frac{1}{2}\gamma,
\end{eqnarray*}
thus validating \eqref{08.16.5}. From this and {\it Step 1}, we have established that
\begin{equation*}
-\iint_Q \widetilde{\varphi}_0 \partial_t \psi \, dx \, dt + \iint_Q (\nabla\widetilde{\varphi}_0 \cdot \nabla\psi) w \, dx \, dt = \iint_Q f \psi \, dx \, dt + \int_\Omega \varphi_0(x) \psi(x,0) \, dx.
\end{equation*}
This confirms that $\widetilde{\varphi}_0$ is a solution of \eqref{12.14.1}.

{\it Step 4}. Given that $\widetilde{\varphi}_0$ and $\widehat{\varphi}_0$ are solutions of \eqref{12.14.1} with initial data $\varphi_0 \in L^2(\Omega)$ and $f w^{-1} \in L^2(\Omega; w)$, by the uniqueness of the solution of \eqref{12.14.1}, we conclude that $\widetilde{\varphi}_0 = \widehat{\varphi}_0$.

{\it Step 5}. Finally, let $\psi \in C^\infty(\overline{Q})$ and $\psi(\cdot, t)$ be compactly supported in $\Omega$ for any $t \in [0,T]$. Multiplying $\psi$ on both sides of \eqref{08.16.1} with $\widehat{\varphi}_k(0) = \varphi_0$ and $f_k = f$, we obtain
\begin{equation*}
\begin{split}
\int_\Omega \widehat{\varphi}_k(T) \psi(T) \, dx - \iint_Q \widehat{\varphi}_k \partial_t \psi \, dx \, dt + \iint_Q w_k \nabla \widehat{\varphi}_k \cdot \nabla\psi \, dx \, dt = \iint_Q f \psi \, dx \, dt + \int_\Omega \widehat{\varphi}_k(0) \psi(0) \, dx.
\end{split}
\end{equation*}
Multiplying $\psi$ on both sides of \eqref{12.14.1}, we obtain
\begin{equation*}
\int_\Omega \widehat{\varphi}_0(T) \psi(T) \, dx - \iint_Q \widehat{\varphi}_0 \partial_t \psi \, dx \, dt + \iint_Q w \nabla \widehat{\varphi}_0 \cdot \nabla\psi \, dx \, dt = \iint_Q f \psi \, dx \, dt + \int_\Omega \widehat{\varphi}_0(0) \psi(0) \, dx.
\end{equation*}
By \eqref{08.16.5} and $\widehat{\varphi}_0 = \widetilde{\varphi}_0$, $\widehat{\varphi}_k(0) = \widehat{\varphi}_0(0) = \varphi_0$ for all $k \in \mathbb{N}$, and $\widehat{\varphi}_k \rightharpoonup \widehat{\varphi}_0$ weakly in $L^2(Q)$, we obtain
\begin{equation*}
\int_\Omega \widehat{\varphi}_k(T) \psi(T) \, dx \to \int_\Omega \widehat{\varphi}_0(T) \psi(T) \, dx.
\end{equation*}
This indicates that $\widehat{\varphi}_k(T) \to \widehat{\varphi}_0(T)$ in $L^2(\Omega)$ since $\psi(T)$ is arbitrary.
\end{proof}

\end{document}
